%!TEX root = iclr2026_conference.tex
% Appendix G (sec:app:related): extended related work.
% Content plan and owner unit: manuscript/paper-contract.md §8.
\section{Extended Related Work}
\label{sec:app:related}

Existing methods reasonably fix predictive horizons and targets to suit their intended objectives. Latent-dynamics control methods fix temporal targets or planning horizons and reach further by recursive one-step rollout with value bootstrapping \citep{hafner2020dreamer,hansen2022tdmpc}. Joint-embedding approaches fix the predictive representation and target at training time \citep{assran2023ijepa,bardes2024vjepa}. Recent latent world models on frozen or end-to-end features inherit the same fixed target \citep{zhou2025dinowm,maes2026leworldmodel}. This task-specific design is a widespread default. We ask whether one action-conditioned model serving heterogeneous decisions should instead select prediction horizon and \rev{description level} jointly at each decision. \rev{On the sealed NovPhy grid, no learned selector ranks launches above fixed $(\Delta=15,\textsf{cont})$ requests on the primary continuous predictor (Section~\ref{sec:res:ranking}).}

Recent work relaxes one dimension of this contract at a time. On the temporal side, variable-length latent world models vary prediction offsets \citep{du2026vlwm}, surprise signals set hierarchical chunk boundaries in open-loop video prediction \citep{iiyama2026sunta}, and reliability estimates adapt rollout depth during imagination \citep{hafez2020improving}. On the description side, published work varies the output form: state-dependent scene resolution \citep{wang2023dynamic} and per-frame abstraction of action forecasts \citep{scarafoni2022islands}. Fixed couplings of the two also exist: hierarchical prediction across fixed time scales \citep{shaj2023multi,villarcorrales2022mspred} and short-detailed versus long-simplified segments in one model-predictive controller \citep{brudigam2021mpc}. \citet{wang2026adajepa} adapt model parameters at test time, not granularity selection. \rev{We use these works as reference points for a narrower interface: one per-query request over horizon and description level, served by a shared action-conditioned predictor.}

Temporal abstraction in reinforcement learning answers a different question. Options and their learned variants decide when and for how long the agent acts \citep{sutton1999options,bacon2017optioncritic}, and \citet{biedenkapp2021temporl} learn when to commit to an action repetition. These methods adapt the agent's control schedule, not what the world model predicts or how far ahead it extends. Our question sits on the prediction side of that boundary.

In the \emph{action-sparse persistent-effect environment}, actions occur infrequently while their physical effects continue autonomously for many fixed simulation steps. A single intervention can set off a cascade through collision onset and active rearrangement before settling. We build on NovPhy \citep{pinto2024novphy}, a published physical-reasoning benchmark \citep{bakhtin2019phyre,bear2021physion,yi2020clevrer} evaluating how agents detect and adapt to novel physical dynamics across scenarios. Each cascade phase, we hypothesize, asks for different prediction granularity: fine temporal resolution and relation-level detail near the intervention, coarse state judgments near settling. We call this tension \emph{granularity mismatch} and treat it as a heterogeneity hypothesis. Effects, if any, depend on the decision's phase, the physical family, and the scored horizon. An average over cells can hide either sign. \rev{The sealed NovPhy split provides no support for exploiting this heterogeneity in launch ranking: among world-model request rankers, a fixed coarse request ranks best (Section~\ref{sec:res:ranking}).}

Granularity can enter a world model on the temporal axis, on the description axis, or on both at once. The paragraphs above survey fixed designs and present them as reasonable task-specific choices. The subsections below cover methods that make at least one axis adaptive, \rev{hierarchical world models,} the closest joint neighbour, and \rev{recent evaluation platforms}. \rev{Each method is described from its own paper. We do not rank the methods against ours.}

\subsection{Temporal adaptation of prediction}
\label{sec:app:rel:temporal}

The temporal-abstraction paragraph above draws the boundary between prediction-side granularity and the control schedule. The options framework, its learned descendants, and TempoRL sit on the acting side of that line, deciding when and for how long the agent acts \citep{sutton1999options,bacon2017optioncritic,biedenkapp2021temporl}. They adapt the agent's commitment rather than the prediction, so we treat them as the boundary of the survey class rather than as members of it.

\begin{revblock}
Per-step temporal adaptation of prediction is prior work. Adaptive skip intervals train a recurrent dynamics model to choose its own prediction intervals \citep{neitz2018adaptive}. \citet{chatterjee2025improving} let the planner optimize the duration of temporally extended actions and select duration ranges with a bandit formulation. In model-based reinforcement learning the learned model then predicts over each chosen duration, which places this work on the prediction side of the boundary above. Both works adapt the temporal axis per step.
\end{revblock}

Within the prediction side, TAWM\rev{, an ICML 2025 paper,} conditions predictions on the time-step size $\Delta t$ and trains across a range of step sizes, so a single model serves tasks observed at different temporal grains \citep{nhu2025tawm}. The grain at inference is imposed by the environment's observation rate rather than chosen per decision, so TAWM tolerates rate changes instead of adapting its grain per state. VLWM predicts future latent states conditioned on action sequences of variable length and extends the horizon through a training curriculum \citep{du2026vlwm}. Its planner is tailored to the variable-length predictor. THICK learns a hierarchy over discrete latent dynamics in which parts of the latent state update sparsely, forming invariant contexts whose durations are learned \citep{gumbsch2024thick}. A higher level predicts when those contexts change. SUNTA varies chunk length at inference inside a hierarchical video predictor, splitting the rollout where surprise spikes, and reports prediction staying accurate to 250 steps where the compared baselines degrade within ten \citep{iiyama2026sunta}. Its scope is open-loop video prediction with a fixed output format. TAWM and VLWM express the prediction in a single latent format, while THICK fixes a two-level discrete latent hierarchy.

\subsection{Description adaptation}
\label{sec:app:rel:description}

The description axis has its own adaptive tradition. VisualPredicator learns abstract world models over an invented first-order predicate language, with an online algorithm that invents predicates and plans over the resulting abstract states \citep{liang2025visualpredicator}. Resolution-adaptive scene representations and per-frame abstraction of action forecasts, both introduced above, also adapt what the model expresses \citep{wang2023dynamic,scarafoni2022islands}. On this axis the adaptation selects the state description. The prediction length remains fixed\rev{, so the horizon is not part of the query}.

\subsection{\texorpdfstring{\rev{Hierarchical world models}}{Hierarchical world models}}
\label{sec:app:rel:hier}

\begin{revblock}
Hierarchical JEPA proposes prediction at several time scales, with short-term, low-level prediction at one level and longer-term, higher-level prediction at the next \citep[\S4.4, \S4.6]{lecun2022path}. In this proposal, longer time scales go with more abstract levels by construction. End-to-end hierarchical JEPA learns such a hierarchy for visual planning \citep{zhang2026hjepa}. HWM learns world models at several temporal scales within one shared latent space and uses the long-horizon model's predictions as subgoals for the short-horizon model \citep{zhang2026hierarchical}. LeJEPA introduced the SIGReg regularizer for joint-embedding training \citep{balestriero2025lejepa}, and LeWorldModel trains a joint-embedding world model end to end from pixels \citep{maes2026leworldmodel}. Our request instead varies horizon and description level independently on one predictor. On the sealed NovPhy grid, tying the two axes, or a constructed three-level hierarchy on our carrier, changes the learned selector's AUC by amounts whose intervals lie within $\pm0.02$ (Appendix~\ref{sec:app:res:controls}). We reproduced end-to-end hierarchical JEPA \citep{zhang2026hjepa} and trained its codebase's re-implementations of HWM and LeWorldModel on NovPhy. Every H-JEPA-family arm we trained failed a pre-declared guard, so these comparisons support no ranking claim in either direction (Appendix~\ref{sec:app:external}).
\end{revblock}

\subsection{Joint state and temporal abstraction}
\label{sec:app:rel:joint}

\rev{SALT jointly learns temporal and state abstractions from offline trajectories and plans with latent skills in a compact latent space \citep{huang2026salt}. Its temporal abstraction is over skills, that is, over behaviours the agent carries out for extended periods. Our request instead sets how many carrier frames one prediction advances after a single launch, and whether predicate feedback enters that prediction, and we score requests against fixed-granularity references.}

\subsection{\texorpdfstring{\rev{Evaluation platforms}}{Evaluation platforms}}
\label{sec:app:rel:eval}

\rev{Two recent platforms evaluate world models.} World-in-World benchmarks world models inside a closed loop with a unified online planning strategy and a standardized action interface \citep{zhang2026worldinworld}. Task success rather than visual fidelity is the primary metric. Reported findings include that visual quality alone does not guarantee task success and that more inference-time compute improves closed-loop performance. WorldBench evaluates video-based world models with physics concepts isolated one at a time, at two levels: intuitive concepts such as object permanence, and low-level constants such as friction coefficients and viscosity \citep{upadhyay2026worldbench}. Both platforms target diagnostic criteria over aggregate quality. \rev{Neither platform treats granularity as the object of evaluation.}

PHYRE certifies task solvability by construction and scores agents by AUCCESS over up to 100 attempts per task~\citep{bakhtin2019phyre}, and I-PHYRE reports human success alongside an experimenter upper bound~\citep{li2024iphyre}. \rev{Physion evaluates event-prediction accuracy, not agent task completion~\citep{bear2021physion}.}

\rev{Several papers report results against horizon or rate:} VLWM reports success against goal offset, TAWM reports success against the evaluation time step, THICK sweeps corridor length, and SUNTA reports horizon-resolved degradation of open-loop video prediction. \rev{Our rollout-error and ranking analyses (Section~\ref{sec:res}) are of the same kind and cover one model family per benchmark.}
